Regulatory frameworks increasingly mandate verifiable accountability for
high-stakes AI decisions, yet adoption of transparency infrastructure
remains negligible. A common objection is cost: operators assume that
cryptographic accountability imposes prohibitive overhead. This paper
tests that assumption empirically.

We construct a total-cost-of-ownership (TCO) model for a three-layer
accountability stack recently proposed in the systems literature:
compile-time evidence binding (BTV, 152~bytes / 1.67~\textmu{}s per
decision), verifiable persistence via transparency logs (BTV-T,
56~bytes / 4.73~ms per decision), and privacy-preserving redaction
via ZK-SNARKs (Accountable Redaction, $\approx$2.1~s / 3.2~kB per
batch of 1{,}000 redactions). Using these measured costs as inputs, we
model the annualised infrastructure expenditure for deployments at
three scales: $10^4$ (regional lender), $10^6$ (national bank), and
$10^8$ (global platform) decisions per year.

We compare this expenditure against the empirical cost of
non-compliance: GDPR enforcement actions totalling~\euro{}4.5B
(2018--2025), SEC off-channel communications fines exceeding~\$2B
(2021--2024), and projected EU~AI~Act penalties of up to 7\% of
global turnover. The analysis reveals a \emph{compliance crossover
point}: the annual decision volume above which verifiable
accountability is cheaper than the expected cost of regulatory
penalties, weighted by historical enforcement probability.

To ground the model beyond cost arithmetic, we conduct semi-structured
interviews with $N=12$ compliance officers and data protection officers
at financial institutions and healthcare providers in Brazil and the EU,
identifying three dominant adoption barriers: (1)~perceived integration
complexity with existing Python/Java pipelines; (2)~uncertainty about
regulatory recognition of compile-time guarantees as ``appropriate
technical measures'' under GDPR Art.~25; and (3)~absence of a neutral
log-operations entity. For each barrier, we map concrete mitigations
already present in the technical stack.

Our findings indicate that for organisations processing more than
approximately $5 \times 10^5$ consequential decisions per year, the
annualised cost of the full accountability stack (\$0.003--\$0.008 per
decision at scale) is below the expected penalty cost under current
enforcement rates. We conclude that the principal obstacle to
verifiable AI accountability is not technical cost but institutional
inertia, and propose a \emph{compliance-credit} mechanism --- analogous
to carbon credits --- whereby early adopters of verified transparency
infrastructure receive regulatory safe-harbour recognition, creating a
market incentive aligned with both economic self-interest and public
accountability.
